\documentclass[../../../include/open-logic-section]{subfiles}

\begin{document}

\olfileid{sth}{replacement}{ref}
\olsection{প্রতিস্থাপন ও প্রতিফলন}

\stagesinex{} দিয়ে প্রতিস্থাপনের ন্যায্যতা
দেওয়ার আমাদের শেষ চেষ্টাটি শুরু হয়
একটি গভীর ও সুন্দর ফল থেকে:\footnote{মনে
করিয়ে দিই: অন্য কথা স্পষ্টভাবে না
বললে সব সূত্রে পরামিতি থাকতে পারে।}
%In this, I will use overlining, such as $x_1, \ldots, x_n$, to
%abbreviate ``$x_1, \ldots, x_n$'':
\begin{thm}[প্রতিফলন ছক]\ollabel{reflectionschema}
যেকোনো সূত্র $\phi$-এর জন্য:
\[
\forall \alpha \exists \beta > \alpha (\forall x_1 \ldots, x_n \in
V_\beta)(\phi(x_1, \ldots, x_n) \liff \phi^{V_\beta}(x_1, \ldots, x_n))
\]
\end{thm}
\noindent
\olref[sth][replacement][strength]{formularelativization}-এর
মতো এখানেও $\phi^{V_\beta}$ হলো
$\phi$-র প্রত্যেক পরিমাণায়ককে
সেট~$V_\beta$-তে সীমাবদ্ধ করে
পাওয়া সূত্র। তাই স্বজ্ঞায় প্রতিফলন
বলে: গোটা স্তরক্রমে $\phi$ সত্য হলে,
স্তরক্রমের যত দূরের খুশি \emph{প্রারম্ভিক
খণ্ডেও} $\phi$ সত্য।
% BN-SRC-795 TARGET-CORRECTION-BEGIN
\emph{উৎস-সংশোধন-নোট।} এটি প্রত্যেক স্থির সূত্রের ছক;
একটি $\beta$-তে সব সূত্র বা গোটা তত্ত্ব একসঙ্গে প্রতিফলিত
হওয়ার দাবি নয়। সূচকগুলি ক্রমসংখ্যা, এবং $\beta$ সূত্র ও
তার নির্দিষ্ট পরামিতির উপর নির্ভর করতে পারে। পরামিতিগুলি
অন্তর্ভুক্ত চাইলে শুরুতে তাদের স্তরাঙ্কের উপরে সূচক নিয়ে
পরে আরও বড় প্রতিফলিত পর্যায় নিই। নিচের উপযুক্ত দুর্বল
রূপে এই পরামিতি-অন্তর্ভুক্তি স্পষ্ট; শুধু $\Z$-তে মোট
$V$-পদ নীরবে যোগ করা হয়নি।
% BN-SRC-795 TARGET-CORRECTION-END

\citet{Montague1961} ও \citet{Levy1960}
দেখিয়েছেন, প্রতিস্থাপন ও প্রতিফলনের
(উপযুক্ত আনুষ্ঠানিক রূপগুলি) $\Z$-এর
সাপেক্ষে তুল্য; ফলে যেকোনো একটি
যোগ করলেই $\ZF$ পাওয়া যায়।
(\olref[sth][replacement][refproofs]{sec}-এ
ফলগুলি প্রমাণ করব।) এই তুল্যতা
থাকায় \stagesinex{} দিয়ে প্রতিফলন
এবং প্রতিস্থাপনকে এইভাবে ন্যায্য
প্রতিপন্ন করার আশা করা যায়:
\stagesinex{} মানলে স্তরক্রম অত্যন্ত
উঁচু হবে—এতটাই যে তার উচ্চতার
সীমা স্থির করতে আমাদের কোনো
বর্ণনাই যথেষ্ট হবে না। একে এভাবে
বোঝা যায়: গোটা স্তরক্রমে কোনো
বাক্য~$\phi$ সত্য হলে স্তরক্রমের
যত দূরের খুশি প্রারম্ভিক খণ্ডেও
তা সত্য। আর এটাই প্রতিফলন।

এটিও প্রতিস্থাপনের অন্তর্নিহিত ন্যায্যতা
দেওয়ার সত্যিই আশাব্যঞ্জক চেষ্টা
বলে মনে হয়। কিন্তু এ বিষয়ে আরও
অনেক কথা বলার আছে; এখানে সে
আলোচনা সম্ভব নয়। এটি সফল কি
না, এবার নিজেরাই বিচার করুন।\footnote{তবে
\citet[95--100]{Incurvati2020} পড়ে
আলোচনা এগিয়ে নিতে পারো।}

\end{document}
